\section*{Hoofdstuk \ref{ch:b3:field-lab-methods-3} --- Veld- en laboratoriummethoden III}
% Values from the Book 5 scratch calculator calc/ch31.py.

\begin{solution}{exo:b3:field-lab-methods-3:1}
(a)~$\qty{0.01}{\milli\gal} = \qty{e-7}{m.s^{-2}}$, en
$\num{e-7}/\num{9.81} = \num{1.0e-8}$: één deel op honderd miljoen.

(b)~$0.01/0.3086 = \qty{0.032}{m}$, ongeveer \qty{3.2}{cm}.

(c)~Samen veranderen de twee effecten de zwaartekracht met
\qty{0.1967}{\milli\gal.m^{-1}}, het verschil $0.3086 - 0.1119$, dus
de hoogte moet met $0.01/0.1967 = \qty{5.1}{cm}$ veranderen.
\end{solution}

\begin{solution}{exo:b3:field-lab-methods-3:2}
$\Delta m = \tfrac43\pi(\qty{50}{m})^3\times\qty{500}{kg.m^{-3}} =
\qty{2.62e8}{kg}$. De piek is $G\Delta m/z^2 =
\num{6.674e-11}\times\num{2.62e8}/\num{e4} = \qty{1.75e-6}{m.s^{-2}}$,
dat is \qty{0.175}{\milli\gal}. De halve breedte is $x_{1/2} = z/1.305 =
\qty{77}{m}$.
\end{solution}

\begin{solution}{exo:b3:field-lab-methods-3:3}
(a)~Met $AM = MN = NB = a$, $K = 2\pi/(1/a - 1/2a - 1/2a + 1/a) =
2\pi a = \qty{62.83}{m}$, en $\rho_a = 62.83\times0.239/0.1 =
\qty{150}{\ohm\metre}$. (b)~$K = \pi(15^2 - 5^2)/(2\times5) = \pi\times
200/10 = \qty{62.83}{m}$. De elektroden liggen dan op $0$, $10$, $20$
en \qty{30}{m}, in de volgorde AMNB: het is de \omterm{def:b3:field-lab-methods-3:resistivity}{Wenner-elektrodeopstelling}
met een elektrodeafstand van \qty{10}{m}.
\end{solution}

\begin{solution}{exo:b3:field-lab-methods-3:4}
(a)~$V = 0.2998/\sqrt5 = \qty{0.134}{m.ns^{-1}}$. (b)~$t =
2\times10/0.134 = \qty{149}{ns}$. (c)~$\delta = 0.0053\sqrt5/\num{e-3} =
\qty{11.9}{m}$, en een bruikbare diepte van ongeveer \qty{36}{m}. (d)~In de klei
$\delta = 0.0053\sqrt{20}/0.01 = \qty{2.4}{m}$, een bruikbare diepte van ongeveer
\qty{7}{m}, en een snelheid van slechts \qty{0.067}{m.ns^{-1}}. De
benadering geldt daar nog steeds: bij \qty{100}{MHz},
$\omega\varepsilon = 2\pi\times\num{e8}\times20\times\num{8.85e-12} =
\qty{0.11}{S.m^{-1}}$, tien keer $\sigma$.
\end{solution}

\begin{solution}{exo:b3:field-lab-methods-3:5}
(a)~De cilinder heeft $\lambda = \pi R^2\Delta\rho$ per lengte-eenheid, en
buiten hem is de aantrekking die van een lijnmassa op de as (stelling van Gauss
voor een coaxiale cilinder): $\Delta g = 2G\pi R^2\Delta\rho\,z/(x^2
+ z^2)$. De helft van de piek wordt bereikt wanneer $x^2 + z^2 = 2z^2$, dus $z =
x_{1/2}$. (b)~De bol ligt op $1.305\times100 = \qty{130}{m}$, de
cilinder op \qty{100}{m}: hetzelfde profiel, als bol gelezen, plaatst het lichaam
\qty{30}{\percent} dieper. De anomalie van een bol neemt sneller af met
de afstand, dus een bol moet dieper liggen om even breed te zijn.
\end{solution}

\begin{solution}{exo:b3:field-lab-methods-3:6}
(a)~Elk punt staat voor $\qty{100}{m}\times\qty{100}{m} =
\qty{e4}{m^2}$, dus $\iint\Delta g\,\mathrm{d}A \approx
\num{-438e-5}\times\num{e4} = \qty{-43.8}{m^3.s^{-2}}$, en $\Delta m
= -43.8/(2\pi\times\num{6.674e-11}) = \qty{-1.04e11}{kg}$. (b)~$1.04/1.74
= \qty{60}{\percent}$. De rest zit in de flanken van de anomalie buiten
de kaart: de graniet reikt tot een diepte van meer dan een kilometer, en de
anomalie van het diepe deel neemt langzaam af met de afstand. Door de kaart
ver genoeg uit te breiden of met een model, is hij terug te winnen.
\end{solution}

\begin{solution}{exo:b3:field-lab-methods-3:7}
(a)~Voor een dunne gang zakt de anomalie tot de helft van haar piek bij $x = \pm z$:
haar volledige breedte op de halve piekhoogte is $2z = \qty{60}{m}$, en drie
metingen erover heen een meetpuntafstand van ten hoogste \qty{20}{m} nodig. (b)~$2\times
220/3 = \qty{147}{m}$. (c)~De benodigde stationsafstand schaalt met de
diepte van het ondiepste lichaam van belang: een meting die op de diepe
gang op \qty{150}{m} is ontworpen, zou de anomalie van de ondiepe gang onderbemonsteren, die
slechts \qty{60}{m} breed is; een meting op het ondiepe lichaam ontworpen ziet ook de
diepe gang, met veel meer metingen dan nodig.
\end{solution}

\begin{solution}{exo:b3:field-lab-methods-3:8}
(a)~$F = 80/8 = 10$ en $\phi = F^{-1/m} = 10^{-1/2} = \num{0.316}$.
(b)~$10^{-1/1.8} = \num{0.278}$ en $10^{-1/2.2} = \num{0.351}$. Een verandering
van \num{0.2} in $m$ verschuift de porositeit met ongeveer vier procentpunten;
de weerstanden, gemeten op enkele procenten nauwkeurig, verschuiven haar veel minder. De
cementatie-exponent bepaalt de onzekerheid, en die moet worden gekalibreerd
op monsters van hetzelfde gesteente.
\end{solution}

\begin{solution}{exo:b3:field-lab-methods-3:9}
(a)~Bij $s = 2h_1$, $(2nh_1/s)^2 = n^2$, dus de termen zijn $k^n(1 +
n^2)^{-3/2}$: $0.3536k$, $0.0894k^2$, $0.0316k^3$. Voor $k = 0.5$:
$\num{0.1768}$, $\num{0.0224}$, $\num{0.0040}$, en $\rho_a/\rho_1 \approx
1 + 2\times0.2031 = \num{1.406}$ (de volledige som geeft \num{1.409}). Voor $k
= -0.5$ wisselen de termen van teken, $-0.1768 + 0.0224 - 0.0040$, en
$\rho_a/\rho_1 \approx \num{0.683}$ (volledige som \num{0.685}). (b)~Voor $s \ll
h_1$ verdwijnt elke term en $\rho_a \to \rho_1$; voor $s \gg h_1$ nadert elke
haakterm naar~1 en $1 + 2\sum k^n = 1 + 2k/(1 - k) = (1 + k)/(1 -
k) = \rho_2/\rho_1$. (c)~$\rho_2/\rho_1 = (1 + k)/(1 - k) = 3$ en $1/3$.
\end{solution}

\begin{solution}{exo:b3:field-lab-methods-3:10}
(a)~In de volgorde AMNB met elektrodeafstand $a$: $AM = a$, $BM = 2a$, $AN = 2a$,
$BN = a$, dus de haakterm is $2/a - 1/a = 1/a$ en $K = 2\pi a$.
(b)~Met A en B op $\mp s$ en M en N op $\mp l$: $AM = BN = s - l$,
$BM = AN = s + l$, de haakterm is $2/(s - l) - 2/(s + l) = 4l/(s^2 -
l^2)$, en $K = \pi(s^2 - l^2)/(2l)$. Wanneer $l \ll s$, $\Delta V_{MN}
\approx 2l\,E$ met $E = 2(-\mathrm{d}V/\mathrm{d}r)$ in het midden, het
veld van beide stroomelektroden; dan is $\rho_a = K\Delta V/I \to
\pi s^2/(2l)\times 4l(-\mathrm{d}V/\mathrm{d}r)/I = 2\pi s^2(-
\mathrm{d}V/\mathrm{d}r)/I$. (c)~De verhouding van de benaderde tot de
exacte factor is $s^2/(s^2 - l^2) = 1/(1 - 0.01) = 1.010$: één procent.
\end{solution}

\begin{solution}{exo:b3:field-lab-methods-3:11}
(a)~Aflezing: \qty{0.010}{\milli\gal}. Hoogte:
$0.1967\times0.05 = \qty{0.010}{\milli\gal}$. Dichtheid: de Bouguer
correctie is $0.04191\,h\,d$ met $d$ in \unit{g.cm^{-3}}, dus
$0.04191\times50\times0.1 = \qty{0.210}{\milli\gal}$. Terrein:
\qty{0.200}{\milli\gal}. Totaal: $\sqrt{0.010^2 + 0.010^2 + 0.210^2 +
0.200^2} = \qty{0.29}{\milli\gal}$. (b)~De dichtheid en de \omterm{def:b3:field-lab-methods-3:terrain}{terreincorrectie}
domineren, elk ongeveer twintig keer de aflezingsfout. Om
\qty{0.1}{\milli\gal} te detecteren, moet de dichtheid van de slab aan
monsters worden gemeten of afgestemd op de gegevens (de waarde die het minst
de topografie weerspiegelt), en de \omterm{def:b3:field-lab-methods-3:terrain}{terreincorrectie} uit een fijn digitaal hoogtemodel.
(c)~$\sqrt{0.010^2 + 0.010^2 + 0.210^2} = \qty{0.21}{\milli\gal}$.
\end{solution}

\begin{solution}{exo:b3:field-lab-methods-3:12}
(a)~De massa die bij de centrale aflezing past is $\Delta m = \Delta g(0)z^2/G
= \num{-4.16e-5}z^2/\num{6.674e-11}$: \qty{-5.61e10}{kg} bij
\qty{300}{m}, \qty{-9.97e10}{kg} bij \qty{400}{m}, \qty{-1.56e11}{kg} bij
\qty{500}{m}. De voorspelde aflezingen bij $x = \pm300$ en $\pm600$ zijn
\num{-1.47} en \num{-0.37}, \num{-2.13} en \num{-0.71}, \num{-2.62}
en \qty{-1.09}{\milli\gal}. De \omterm{def:b3:field-lab-methods-3:inversion-practice}{misfits} zijn $\chi^2 = 1238$, $3.7$ en
$856$. (b)~Alleen het model bij \qty{400}{m}, met $\chi^2 = 3.7$ voor vijf
gegevens, past binnen de fouten. (c)~Ja: elk bolsymmetrisch lichaam
met dezelfde massa en hetzelfde middelpunt geeft exact dezelfde anomalie, en andere
vormen kunnen binnen de fouten vallen. Een dichtheid gemeten aan monsters, in een
boorgat of met een andere methode zou beslissen.
\end{solution}

\begin{solution}{pb:b3:field-lab-methods-3:1}
\textbf{1.} $(2846.708 - 2846.512)/\qty{2.0}{h} = \qty{0.098}{\milli\gal.h^{-1}}$.

\textbf{2.} $2851.075 - 0.098\times0.5 - 2846.512 = \qty{4.514}{\milli\gal}$;
$2850.497 - 0.088 - 2846.512 = \qty{3.897}{\milli\gal}$;
$2849.392 - 0.127 - 2846.512 = \qty{2.753}{\milli\gal}$.

\textbf{3.} De hoogten boven de basis zijn $-27.1$, $-30.6$ en
\qty{-18.1}{m}; door $0.1967\,\Delta h$ op te tellen volgen de Bouguer-anomalieën
$\num{-0.817}$, $\num{-2.122}$ en \qty{-0.808}{\milli\gal}. De
meetpunten liggen lager dan de basis en dichter bij het middelpunt van de Aarde,
dus hun ruwe aflezingen zijn hoger; de correctie verwijdert dat en
laat een laagte boven het graniet over.

\textbf{4.} Het extra $0.098 - 0.040 = \qty{0.058}{\milli\gal.h^{-1}}$,
\qty{0.12}{\milli\gal} over de lus, is het getij van de vaste Aarde. Over twee
uur verandert het bijna lineair, dus een rechte lijn door de twee
basisaflezingen verwijdert het samen met de drift.

\textbf{5.} De fouten zijn $\num{0.008}$, $\num{0.009}$ en
\qty{0.017}{\milli\gal}. Door elke hoogte af te ronden op \qty{0.1}{m} kan
een hoogteverschil tot \qty{0.1}{m} onjuist zijn, dat is
\qty{0.020}{\milli\gal}: de afronding verklaart de fouten, het getij dat
de rechte lijn overlaat, is kleiner dan \qty{0.001}{\milli\gal}.

\textbf{6.} $0.01/0.1967 = \qty{5.1}{cm}$. De \qty{3.3}{cm} uit de handleidingen is
$0.01/0.3086$: de vrije-luchtanomalie hangt niet af van de dichtheid
die voor de Bouguer-correctie is gekozen, en de hoogte moet ook daarvoor dienen.

\textbf{7.} $x_{1/2} = (3.952 - 3.048)/2 = \qty{0.452}{km}$, en $z = 1.305
\times452 = \qty{590}{m}$ onder het maaiveld.

\textbf{8.} $\Delta m = \Delta g\,z^2/G = \num{-2.293e-5}\times590^2/
\num{6.674e-11} = \qty{-1.20e11}{kg}$. Met $\Delta\rho =
\qty{-130}{kg.m^{-3}}$ is het volume \qty{9.2e8}{m^3} en de straal
\qty{603}{m}, groter dan de diepte van het middelpunt: de bol zou uit de
grond steken. Geen bol met het dichtheidsverschil van het graniet
past; het lichaam moet tot dicht bij het oppervlak reiken.

\textbf{9.} Het graniet loopt van het oppervlak tot een grote diepte. Het
ondiepe deel, nabij de meetpunten, maakt de piek smal, en de regel, die
alleen de breedte kent, plaatst één bol op een matige diepte; het
diepe deel voegt brede flanken toe die de bol ontbreken (de figuur van het
profiel), en de massa die zij bevatten wordt gemist.

\textbf{10.} $2\pi G\Delta\rho = 2\pi\times\num{6.674e-11}\times(-130) =
\qty{-5.45e-8}{s^{-2}}$. Met $z_1 = 0$, $z_2 = \qty{1645}{m}$, $a =
\qty{461}{m}$: de haakjesterm is $1645 + 461 - 1708.4 = \qty{397.6}{m}$,
en $\Delta g = \qty{-2.17}{\milli\gal}$. Bij een onbeperkte diepte nadert de
haakjesterm $a$, en $\Delta g = \qty{-2.51}{\milli\gal}$.

\textbf{11.} De haakjesterm moet $2.29/\num{5.45e-3} = \qty{420}{m}$ zijn, wat een
ondergrens op ongeveer \qty{2.6}{km} geeft. Elke ondergrens tussen ongeveer
\qty{1.6}{km} en een onbeperkte diepte verandert de piek met slechts enkele tienden
van een milligal, en een iets bredere of dichtere massa geeft dezelfde piek met
een ondiepere ondergrens: de ondergrens is niet op te lossen.

\textbf{12.} Met $z_1 = \qty{6}{m}$: $1639 + 461.0 - 1708.4 = 391.6$, en
$\Delta g = \qty{-2.14}{\milli\gal}$, \qty{0.03}{\milli\gal} minder dan
met de bovenkant aan het maaiveld. Dat ligt onder de terrein- en
dichtheidsfouten van het onderzoek: de zwaartekracht kan niet zeggen of het graniet
tot aan het oppervlak onder het alluvium reikt. De sondering wel.

\textbf{13.} $K = \pi(20^2 - 1^2)/(2\times1) = \qty{626.7}{m}$, en $\Delta V =
\rho_a I/K = 821\times0.050/626.7 = \qty{65.5}{mV}$.

\textbf{14.} $V = 2\times1/\sqrt{17.1^2 - 10.7^2} = 2/13.34 =
\qty{0.150}{m.ns^{-1}}$, $\varepsilon_r = (0.2998/0.150)^2 = 4.0$, een
droog zand; de pijp ligt $0.150\times10.7/2 = \qty{0.80}{m}$ diep.

\textbf{15.} $0.150\times20.0/2 = \qty{1.50}{m}$.

\textbf{16.} $20/821 = \num{0.0244}$, $30/1309 = \num{0.0229}$, $40/1684 =
\num{0.0238}$: $S \approx \qty{0.0237}{S}$.

\textbf{17.} De droge laag draagt $1.5/600 = \qty{0.0025}{S}$ bij, de
verzadigde $\qty{0.0212}{S}$, dus de dikte ervan is $0.0212\times222 =
\qty{4.7}{m}$, en het graniet ligt $1.5 + 4.7 = \qty{6.2}{m}$ onder de
vloer. De schatting is te hoog: het graniet is niet oneindig
weerstandig, dus bij deze afstanden ligt $\rho_a$ nog onder de lijn
$s/S$, en $s/\rho_a$ overschat $S$. (Het model van de regio heeft
\qty{6}{m} alluvium.)

\textbf{18.} $\phi = (20/222)^{1/2} = \num{0.30}$. Met $m = 1.3$, $\phi =
0.0901^{1/1.3} = \num{0.16}$; met $m = 2.8$, $\phi = \num{0.42}$.

\textbf{19.} Door de basis van de lijn af te trekken volgt, van west naar oost:
3.4, 4.1, 4.8, 5.9, 7.4, 9.2, 11.5, 14.2, 17.2, 19.8, 21.0, 20.1, 17.6,
14.5, 11.5, 9.1, 7.2, 5.7, 4.6, 3.8 en \qty{3.1}{nT}.

\textbf{20.} Ongecorrigeerd is de piek \qty{28.0}{nT} bij \qty{2.20}{km}, en
de helft daarvan, \qty{14.0}{nT}, wordt naar het oosten nooit bereikt: de lijn
eindigt bij \qty{15.2}{nT}, omdat het veld tijdens de
wandeling met \qty{12}{nT} steeg.

\textbf{21.} De piek is \qty{21.0}{nT} bij \qty{2.20}{km}; de helft ervan is
\qty{10.5}{nT}. West: tussen \num{9.2} bij \num{1.95} en \num{11.5} bij
\qty{2.00}{km}, bij $1.95 + 0.05\times1.3/2.3 = \qty{1.978}{km}$. Oost:
tussen \num{11.5} bij \num{2.40} en \num{9.1} bij \qty{2.45}{km}, bij
$2.40 + 0.05\times1.0/2.4 = \qty{2.421}{km}$. De halve breedte is $(2.421
- 1.978)/2 = \qty{0.221}{km}$.

\textbf{22.} Een symmetrische anomalie in een veld met een inclinatie van \ang{65} betekent dat
alleen de verticale magnetisatie werkt: het lichaam heeft zijn strekking langs de magnetische
meridiaan, noord--zuid. De regel van
\cref{prop:b3:field-lab-methods-3:dyke} gaat uit van een dunne, verticale,
tweedimensionale gang die ver naar beneden reikt, door inductie gemagnetiseerd,
waarvan de bovenkant werkt als een lijn van magnetische lading.

\textbf{23.} $\kappa t = 2\pi z\,\Delta T_\mathrm{max}/(F\sin^2 I) = 2\pi\times
221\times21.0/(48000\times0.8214) = \qty{0.74}{m}$. Met $\kappa =
0.03$, $t = \qty{25}{m}$. Alleen het product ligt vast: een gang die twee keer zo
dik is maar de halve susceptibiliteit heeft, geeft dezelfde anomalie.

\textbf{24.} De anomalie is \qty{443}{m} breed op de helft van haar piek en de meetpunten
liggen \qty{50}{m} uit elkaar: negen aflezingen erover, drie keer het
minimum van de regel.

\textbf{25.} De bovenkant van de gang ligt
\textbf{\qty[text-series-to-math=true]{220}{m}} onder het maaiveld. Het
geologische model plaatst de discordantie op dit punt \qty{222}{m} diep:
de gang stijgt op tot de basis van de bedekking en wordt daar afgesneden, en is
ouder dan de Trias-zandsteen die hem bedekt.
\end{solution}
